% ─────────────────────────────────────────────────────────────────────────────
% natbib backend test
%
% Verifies citemsp works with natbib (the path used by revtex4-2, aastex,
% elsarticle, etc.).  Uses plain article + natbib to avoid needing the
% journal class files installed locally.
%
% Build:  pdflatex test-natbib && bibtex test-natbib && pdflatex test-natbib && pdflatex test-natbib
%    or:  latexmk -pdf test-natbib.tex
% ─────────────────────────────────────────────────────────────────────────────
\documentclass[11pt, a4paper]{article}

\usepackage{lmodern}
\usepackage[margin=2.5cm]{geometry}
\usepackage[numbers, sort&compress]{natbib}

\usepackage{citemsp}
\usepackage[colorlinks=true, citecolor=blue]{hyperref}

\begin{document}

\section*{citemsp --- natbib Backend Test}

% ═════════════════════════════════════════════════════════════════════════════
\subsection*{1. Plain citations}
% ═════════════════════════════════════════════════════════════════════════════

\noindent
\texttt{\textbackslash cite} baseline: \cite{ref1,ref2,ref3,ref4,ref5} \\
\texttt{\textbackslash citemsp} same keys: \citemsp{ref1, ref2, ref3, ref4, ref5} \\
Single: \citemsp{ref3}

% ═════════════════════════════════════════════════════════════════════════════
\subsection*{2. Default locators: section and paragraph}
% ═════════════════════════════════════════════════════════════════════════════

\noindent
Section only:              \citemsp{ref3/6.1} \\
Section + paragraph:       \citemsp{ref3/6.1/2} \\
Paragraph only (sub slot): \citemsp{ref3//5} \\
Deep subsection + para:    \citemsp{ref3/3.2.1/4}

% ═════════════════════════════════════════════════════════════════════════════
\subsection*{3. Built-in prefix locators}
% ═════════════════════════════════════════════════════════════════════════════

\noindent
\begin{tabular}{llll}
\textbf{Prefix} & \textbf{Type} & \textbf{In super slot} & \textbf{In sub slot} \\[2pt]
\texttt{d} (definition) & $\triangleq$  & \citemsp{ref3/d5}        & \citemsp{ref3/3.2/d5} \\
\texttt{e} (equation)   & eq.           & \citemsp{ref3/e4.3}      & \citemsp{ref3/3.2/e4.3} \\
\texttt{f} (figure)     & fig.          & \citemsp{ref3/f3}        & \citemsp{ref3/6.1/f3} \\
\texttt{n} (footnote)   & $\dagger$     & \citemsp{ref4/n7}        & \citemsp{ref4/1/n7} \\
\texttt{p} (page)       & pp.           & \citemsp{ref3/p42}       & \citemsp{ref3/3.2/p42} \\
\texttt{t} (table)      & $\boxplus$    & \citemsp{ref5/t2}        & \citemsp{ref5/8.4/t2} \\
\end{tabular}

% ═════════════════════════════════════════════════════════════════════════════
\subsection*{4. Mixed locator types}
% ═════════════════════════════════════════════════════════════════════════════

\noindent
\verb|\citemsp{ref3/6.1/2, ref5/p520, ref4/1/e4.3}|\\
Output: \citemsp{ref3/6.1/2, ref5/p520, ref4/1/e4.3}

% ═════════════════════════════════════════════════════════════════════════════
\subsection*{5. Italic context}
% ═════════════════════════════════════════════════════════════════════════════

\noindent
Upright: The field equations~\citemsp{ref3/3.2/e4.3} can be derived.

\medskip\noindent
\textit{Italic: Hawking radiation~\citemsp{ref4/1/2} implies black holes radiate.}

% ═════════════════════════════════════════════════════════════════════════════
\subsection*{6. Realistic paragraph}
% ═════════════════════════════════════════════════════════════════════════════

The Schwarzschild solution~\citemsp{ref2/1/3} was found shortly after the
field equations~\citemsp{ref1} were published. A modern treatment can be
found in standard textbooks~\citemsp{ref3/6.1/2, ref5/p520}. The
singularity theorems~\citemsp{ref6, ref4/1/e4.3} later showed that
gravitational collapse is generic.

\bibliographystyle{unsrtnat}
\bibliography{sandbox-refs}

\end{document}
